% GENERATED FILE. Edit canonical lessons, then run npm run generate:books.
% canonical-source-hash: fnv1a64:257733a8b5ccd51d
% canonical-lessons: TA-C466-tuniccal, TA-C466-kolai, TA-C466-ancu, TA-C466-patunku, TA-C466-mirattu

\chapter{Courage and Cowardice}
\label{ch:ta-courage-and-cowardice}
\hlchaptermodality{466}

\begin{chapteropening}
\textbf{By the end of this chapter:} \emph{I can talk about courage, cowards and threats.}

Complete the last lesson of chapter 466: I can talk about courage, cowards and threats.
\end{chapteropening}

\section[\texorpdfstring{\ta{துணிச்சல்} (tuṇiccal)}{துணிச்சல் (tuṇiccal)}]{\ta{துணிச்சல்} (tuṇiccal) — courage; daring}

\label{lesson:TA-C466-tuniccal}



\begin{quote}
Before the new one: say the Tamil for infatuation, then the Tamil for to be moved.
\end{quote}

\subsection*{You'll want to know: \ta{துணிச்சல்}}
\textbf{\ta{துணிச்சல்}} — \emph{tuṇiccal} — \textquotedblleft{}courage; daring\textquotedblright{}.

\begin{quote}
\ta{இரவில்} \ta{தனியாகப்} \ta{போக} \ta{அவளுக்குத்} \ta{துணிச்சல்} \ta{இருந்தது}.
\end{quote}

\begin{quote}
\emph{iravil taṉiyākap pōka avaḷukkut tuṇiccal iruntatu.}
\end{quote}

\begin{quote}
She had the courage to go alone at night.
\end{quote}

From \ta{துணி} \textquotedblleft{}to dare\textquotedblright{}, not \ta{துணி} \textquotedblleft{}cloth\textquotedblright{}.

\begin{cousinweb}
\textbf{\ta{துணிச்சல்}} (\emph{tuṇiccal}) is native Dravidian, a noun from \ta{துணி} \emph{tuṇi}, \textquotedblleft{}to dare\textquotedblright{} (DEDR 3306), which lists it.
\end{cousinweb}

\begin{grammarlens}[title={What you've built}]
One more word for describing what you lived and felt.
\end{grammarlens}

\subsection*{Your turn}
\begin{itemize}
\raggedright
  \item \emph{Say it:} \emph{tuṇiccal}
  \item \emph{Say it:} \emph{tuṇiccal}, once more
  \item \emph{Say it:} \emph{iravil taṉiyākap pōka avaḷukkut tuṇiccal iruntatu.} — She had the courage to go alone at night.
  \item \emph{Recall:} say the Tamil for infatuation, then the Tamil for to be moved, then say \emph{tuṇiccal} again
\end{itemize}

\begin{tcolorbox}[breakable,colback=teal!4,colframe=teal!35!black,title={Before you move on}]
What does \ta{துணிச்சல்} mean? (\textquotedblleft{}Courage; daring\textquotedblright{}.) Say it once more.
\end{tcolorbox}
\section[\texorpdfstring{\ta{கோழை} (kōḻai)}{கோழை (kōḻai)}]{\ta{கோழை} (kōḻai) — a coward}

\label{lesson:TA-C466-kolai}



\begin{quote}
Before the new one: say the Tamil for to be moved, then the Tamil for courage.
\end{quote}

\subsection*{You'll want to know: \ta{கோழை}}
\textbf{\ta{கோழை}} — \emph{kōḻai} — \textquotedblleft{}a coward\textquotedblright{}.

\begin{quote}
\ta{பயந்து} \ta{ஓடியதால்} \ta{அவனைக்} \ta{கோழை} \ta{என்று} \ta{சொன்னார்கள்}.
\end{quote}

\begin{quote}
\emph{payantu ōṭiyatāl avaṉaik kōḻai eṉṟu coṉṉārkaḷ.}
\end{quote}

\begin{quote}
Because he ran away in fear, they called him a coward.
\end{quote}

X-\ta{ஐ} Y \ta{என்று} \ta{சொல்}, to call X a Y.

\begin{cousinweb}
\textbf{\ta{கோழை}} (\emph{kōḻai}) is native Dravidian (DEDR 2250), where it names both timidity and a timid person.
\end{cousinweb}

\begin{grammarlens}[title={What you've built}]
One more word for describing what you lived and felt.
\end{grammarlens}

\subsection*{Your turn}
\begin{itemize}
\raggedright
  \item \emph{Say it:} \emph{kōḻai}
  \item \emph{Say it:} \emph{kōḻai}, once more
  \item \emph{Say it:} \emph{payantu ōṭiyatāl avaṉaik kōḻai eṉṟu coṉṉārkaḷ.} — Because he ran away in fear, they called him a coward.
  \item \emph{Recall:} say the Tamil for to be moved, then the Tamil for courage, then say \emph{kōḻai} again
\end{itemize}

\begin{tcolorbox}[breakable,colback=teal!4,colframe=teal!35!black,title={Before you move on}]
What does \ta{கோழை} mean? (\textquotedblleft{}A coward\textquotedblright{}.) Say it once more.
\end{tcolorbox}
\section[\texorpdfstring{\ta{அஞ்சு} (añcu)}{அஞ்சு (añcu)}]{\ta{அஞ்சு} (añcu) — to dread; to be afraid}

\label{lesson:TA-C466-ancu}



\begin{quote}
Before the new one: say the Tamil for courage, then the Tamil for a coward.
\end{quote}

\subsection*{You'll want to know: \ta{அஞ்சு}}
\textbf{\ta{அஞ்சு}} — \emph{añcu} — \textquotedblleft{}to dread; to be afraid\textquotedblright{}.

\begin{quote}
\ta{அவன்} \ta{யாருக்கும்} \ta{அஞ்சவில்லை}.
\end{quote}

\begin{quote}
\emph{avaṉ yārukkum añcavillai.}
\end{quote}

\begin{quote}
He was afraid of no one.
\end{quote}

Past (I …): \ta{அஞ்சினேன்}, \emph{añciṉēṉ}.

More literary than \ta{பயப்படு}, the everyday verb. What is feared takes -\ta{க்கு}.

\begin{cousinweb}
\textbf{\ta{அஞ்சு}} (\emph{añcu}) is native Dravidian (DEDR 55); the same entry gives \ta{அச்சம்} \emph{accam}, \textquotedblleft{}fear\textquotedblright{}.
\end{cousinweb}

\begin{grammarlens}[title={What you've built}]
One more word for describing what you lived and felt.
\end{grammarlens}

\subsection*{Your turn}
\begin{itemize}
\raggedright
  \item \emph{Say it:} \emph{añcu}
  \item \emph{Say it:} \emph{añcu}, once more
  \item \emph{Say it:} \emph{avaṉ yārukkum añcavillai.} — He was afraid of no one.
  \item \emph{Recall:} say the Tamil for courage, then the Tamil for a coward, then say \emph{añcu} again
\end{itemize}

\begin{tcolorbox}[breakable,colback=teal!4,colframe=teal!35!black,title={Before you move on}]
What does \ta{அஞ்சு} mean? (\textquotedblleft{}To dread; to be afraid\textquotedblright{}.) Say it once more.
\end{tcolorbox}
\section[\texorpdfstring{\ta{பதுங்கு} (patuṅku)}{பதுங்கு (patuṅku)}]{\ta{பதுங்கு} (patuṅku) — to cower; to crouch}

\label{lesson:TA-C466-patunku}



\begin{quote}
Before the new one: say the Tamil for a coward, then the Tamil for to dread.
\end{quote}

\subsection*{You'll want to know: \ta{பதுங்கு}}
\textbf{\ta{பதுங்கு}} — \emph{patuṅku} — \textquotedblleft{}to cower; to crouch\textquotedblright{}.

\begin{quote}
\ta{பூனை} \ta{சத்தம்} \ta{கேட்டுக்} \ta{கட்டிலுக்குக்} \ta{கீழே} \ta{பதுங்கியது}.
\end{quote}

\begin{quote}
\emph{pūṉai cattam kēṭṭuk kaṭṭilukkuk kīḻē patuṅkiyatu.}
\end{quote}

\begin{quote}
Hearing the noise, the cat crouched under the bed.
\end{quote}

Past (I …): \ta{பதுங்கினேன்}, \emph{patuṅkiṉēṉ}.

To hide by crouching low, as an animal does.

\begin{cousinweb}
\textbf{\ta{பதுங்கு}} (\emph{patuṅku}) is native Dravidian (DEDR 3912); its strong partner \ta{பதுக்கு} \emph{patukku} means to hide something away.
\end{cousinweb}

\begin{grammarlens}[title={What you've built}]
One more word for describing what you lived and felt.
\end{grammarlens}

\subsection*{Your turn}
\begin{itemize}
\raggedright
  \item \emph{Say it:} \emph{patuṅku}
  \item \emph{Say it:} \emph{patuṅku}, once more
  \item \emph{Say it:} \emph{pūṉai cattam kēṭṭuk kaṭṭilukkuk kīḻē patuṅkiyatu.} — Hearing the noise, the cat crouched under the bed.
  \item \emph{Recall:} say the Tamil for a coward, then the Tamil for to dread, then say \emph{patuṅku} again
\end{itemize}

\begin{tcolorbox}[breakable,colback=teal!4,colframe=teal!35!black,title={Before you move on}]
What does \ta{பதுங்கு} mean? (\textquotedblleft{}To cower; to crouch\textquotedblright{}.) Say it once more.
\end{tcolorbox}
\section[\texorpdfstring{\ta{மிரட்டு} (miraṭṭu)}{மிரட்டு (miraṭṭu)}]{\ta{மிரட்டு} (miraṭṭu) — to threaten; to intimidate}

\label{lesson:TA-C466-mirattu}



\begin{quote}
Before the new one: say the Tamil for to dread, then the Tamil for to cower.
\end{quote}

\subsection*{You'll want to know: \ta{மிரட்டு}}
\textbf{\ta{மிரட்டு}} — \emph{miraṭṭu} — \textquotedblleft{}to threaten; to intimidate\textquotedblright{}.

\begin{quote}
\ta{பெரிய} \ta{பையன்கள்} \ta{அவனை} \ta{மிரட்டினார்கள்}.
\end{quote}

\begin{quote}
\emph{periya paiyaṉkaḷ avaṉai miraṭṭiṉārkaḷ.}
\end{quote}

\begin{quote}
The big boys intimidated him.
\end{quote}

Past (I …): \ta{மிரட்டினேன்}, \emph{miraṭṭiṉēṉ}.

The one threatened takes -\ta{ஐ}.

\begin{cousinweb}
\textbf{\ta{மிரட்டு}} (\emph{miraṭṭu}) is native Dravidian (DEDR 5075), the causative of \ta{மிரள்} \emph{miraḷ}, \textquotedblleft{}to be frightened\textquotedblright{}.
\end{cousinweb}

\begin{grammarlens}[title={What you've built}]
One more word for describing what you lived and felt.
\end{grammarlens}

\subsection*{Your turn}
\begin{itemize}
\raggedright
  \item \emph{Say it:} \emph{miraṭṭu}
  \item \emph{Say it:} \emph{miraṭṭu}, once more
  \item \emph{Say it:} \emph{periya paiyaṉkaḷ avaṉai miraṭṭiṉārkaḷ.} — The big boys intimidated him.
  \item \emph{Recall:} say the Tamil for to dread, then the Tamil for to cower, then say \emph{miraṭṭu} again
\end{itemize}

\begin{tcolorbox}[breakable,colback=teal!4,colframe=teal!35!black,title={Before you move on}]
What does \ta{மிரட்டு} mean? (\textquotedblleft{}To threaten; to intimidate\textquotedblright{}.) Say it once more.
\end{tcolorbox}
